% What may stand between an alignment that filled a display and the $$ that
% closes it.
%
% Measured against pdftex: only $$, assignments, and \relax. Anything else
% -- a letter, \hbox, \vskip, \par, \eqno -- gets "Missing $$ inserted",
% the $$ goes in, the display closes, and the command is read again in
% whatever mode that leaves. \eqno then draws a second diagnostic, "You
% can't use `\eqno' in horizontal mode", which is how trip line 254 gets
% the two error messages its own comment promises.
%
% \setbox is the exception among assignments: it is refused outright with
% "Improper \setbox" (see a-font-names-itself-before-its-metrics).
%
% Blanks stand between those assignments -- an alignment usually ends a
% line -- and the reference skips them rather than closing the display on
% one.
%
% Deciding this needs a predicate for "does this command assign", built
% from the main loop's own cases and then checked command by command
% against the reference: it admits the prefixes, \parshape, \textfont,
% \patterns and the interaction modes along with the obvious ones, and
% refuses \message, \mark, \penalty, \unskip, \showthe, \lowercase,
% \ignorespaces, \setlanguage, \afterassignment and \aftergroup.
\tracingonline=1
\font\tt=cmr10 \tt
$$\halign{#\cr A\cr}\count0=1 $$\message{[assignment is allowed]}\par
$$\halign{#\cr A\cr}\relax$$\message{[relax is allowed]}\par
$$\halign{#\cr A\cr}$$\message{[a bare close is silent]}\par
$$\halign{#\cr A\cr}\eqno x$$\message{[eqno draws two]}\par
$$\halign{#\cr A\cr}
\count0=1 $$\message{[a blank line is skipped]}\par
\end
